% Catatan: berkas ini dipindahkan dari Bab IV (chapter-4.tex) hasil restrukturisasi,
% karena isinya adalah persiapan alat/lingkungan pengembangan yang secara metodologis
% lebih tepat berada di Bab III (Metode Penelitian), bukan Bab IV (Hasil dan Pembahasan).
% Konten belum disisipkan ke chapter-3.tex -- simpan di sini dulu sebagai draf siap pakai.
% Seluruh isi berikut tetap valid dan terverifikasi langsung di lingkungan kerja.

\subsection{Persiapan Lingkungan Pengembangan}

Lingkungan pengembangan dibangun di atas Linux Ubuntu 24.04 LTS. Komponen \textit{toolchain} yang dipasang meliputi Yosys (sintesis RTL), nextpnr-xilinx (\textit{place-and-route}), openXC7 (paket \texttt{fasm2frames}/\texttt{xc7frames2bit}/\texttt{bit2fasm} untuk pembuatan \textit{bitstream}), riscv64-unknown-elf-gcc v13.2.0 (\textit{cross-compiler firmware}), dan Icarus Verilog (\texttt{iverilog}) untuk simulasi fungsional sebelum sintesis penuh dijalankan.

Paket \texttt{openxc7} tidak tersedia di Snap Store publik, sehingga tidak dapat dipasang lewat \texttt{sudo snap install openxc7} secara langsung. Sebagai gantinya, berkas \texttt{.snap} diunduh langsung dari GitHub Releases resmi proyek (\texttt{openXC7/yosys-snap} dan \texttt{openXC7/openXC7-snap}) dan dipasang secara lokal menggunakan \texttt{snap install --classic --dangerous}. \texttt{pypy3} dipasang terpisah lewat \texttt{apt}, karena dibutuhkan oleh \texttt{bbaexport.py} untuk membangkitkan \textit{chip database} (\texttt{chipdb}) namun tidak disertakan dalam paket \texttt{.snap}. \texttt{Makefile} proyek juga mengasumsikan struktur direktori \textit{sibling}, yaitu \texttt{openXC7.mk} dan \texttt{chipdb/} berada satu level di atas \textit{repository} kerja, bukan di dalamnya.

Selama tahap ini, ditemukan lima hambatan instalasi yang terverifikasi langsung di lingkungan kerja:

\begin{itemize}
	\item \textbf{\texttt{make synth} gagal mencari \texttt{openXC7.mk}} -- penyebabnya struktur direktori \textit{sibling} belum tersedia; solusinya membuat skrip \texttt{install.sh} yang menyalin \texttt{openXC7.mk} dan menyiapkan direktori \texttt{chipdb/} kosong satu level di atas \textit{repository} secara otomatis.
	\item \textbf{\texttt{bbasm} gagal menulis \texttt{chipdb}} -- penyebabnya variabel \texttt{CHIPDB} didefinisikan sebagai \textit{path} relatif; solusinya mengubahnya menjadi \textit{path} absolut (\texttt{CHIPDB = \$(abspath ../chipdb)}).
	\item \textbf{\texttt{snap install openxc7} gagal} -- penyebabnya paket tidak terdaftar di Snap Store publik; solusinya mengunduh berkas \texttt{.snap} langsung dari GitHub Releases proyek \texttt{openXC7}, sebagaimana diuraikan di atas.
	\item \textbf{\texttt{pypy3: No such file or directory}} -- penyebabnya \texttt{pypy3} bukan bagian dari paket \texttt{.snap}; solusinya memasangnya terpisah lewat \texttt{apt install pypy3}.
	\item \textbf{\texttt{main.c} gagal dikompilasi} -- penyebabnya tipe \texttt{uint8\_t} dipakai tanpa \texttt{\#include <stdint.h>}; solusinya menambahkan berkas \textit{header} tersebut.
\end{itemize}
